\begin{recipeDP}
    [
        preparationtime = {\SI{45}{\minute}},
        portion = {4 Portionen},
        source = {elavegan.com}
    ]
    {Gulasch}

    \graph
        {
            small=Hauptgerichte/Suppen_und_Eintoepfe/Gulasch/small.jpg,
            big=Hauptgerichte/Suppen_und_Eintoepfe/Gulasch/big.jpg
        }

    % \introduction{
    %     Ein herzhafter, aromatischer Eintopf mit Paprika, Kartoffeln, Zwiebeln und Sojaschnetzeln. Perfekt mit Kartoffelpüree, Brot oder Nudeln.
    % }

    \ingredients{
        2 & Zwiebeln \\
        2 & rote Paprika\\
        2 & grüne Paprika \\
        3 & \addtoidx{Kartoffeln} \\
        2 & Karotten \\
        3 & Tomaten \\
        5 & Knoblauchzehen \\
        \SI{4}{\EL} & Paprikapulver, edelsüß \\
        \SI[parse-numbers = false]{\nicefrac{1}{2}}{\TL} & Kreuzkümmel \\
        \SI{1}{\TL} & Salz \\
        \SI{1}{\TL} & Zwiebelpulver \\
        \SI[parse-numbers = false]{\nicefrac{1}{4}}{\TL} & Räucherpaprika \\
        & Cayennepfeffer\\
        & Pfeffer \\
        \SI{2}{\TL} & Zucker \\
        2 & Lorbeerblätter \\
        \SI{480}{\g} & Passata \\
        \SI{480}{\ml} & Gemüsebrühe \\
        \SI{130}{\g} & \addtoidx{Sojaschnetzel} \\
        \SI{300}{\ml} & Wasser \\
        \SI{2}{\EL} & Sojasauce \\
        \SI[parse-numbers = false]{\nicefrac{1}{2}}{\EL} & Balsamico-Essig \\
        \SI{80}{\ml} & Kokosmilch \\
        & frische Kräuter
    }

    \preparation{
        \step Die Sojaschnetzel in \SI{300}{\ml} Wasser für etwa \SI{10}{\minute} einweichen, dann abtropfen lassen und gut ausdrücken.

        \step In der Zwischenzeit Zwiebeln, Paprikaschoten, Kartoffeln, Karotten und Tomaten klein schneiden. Knoblauch fein hacken.
        Einen großen Topf erhitzen und die Zwiebeln zusammen mit den Paprikaschoten für etwa\SI{4}{\minute} bei mittlerer Hitze dünsten.

        \step Karotten, Knoblauch, Kartoffeln, Tomaten, alle Gewürze, Tomatensauce, Lorbeerblätter und die Gemüsebrühe oder das Wasser hinzufügen. Alles gut verrühren und zum Kochen bringen.

        \step Das Gulasch etwa \SI{15}{\minute} köcheln lassen. Dann die eingeweichten Sojaschnetzel, die Sojasauce und den Balsamico-Essig unterrühren.

        \step Weitere \SI{20}{\minute} kochen, oder bis Kartoffeln und Gemüse weich sind. Zum Schluss die Kokosmilch unterrühren und mit Salz, Pfeffer und eventuell etwas Chilipulver abschmecken.

        \step Mit Kartoffelpüree, Brot oder Nudeln servieren. Mit frischen Kräutern garnieren.
    }

    % \hint{
        
    % }

\end{recipeDP}
